\section{Conclusiones}\label{sec:conclusiones}

La introducción planteó una pregunta operativa: con una capacidad de inspección limitada, a
qué clientes visitar primero. El historial de facturación permite responderla mejor que el
azar y mejor que una regla basada en cuánto historial tiene cada cliente. En una campaña que
visita al 1\,\% de los clientes, el modelo propuesto encuentra 118 fraudes donde la selección
al azar encontraría 15 y la regla de exposición 57, y lo hace sobre clientes que no participaron
en ninguna decisión del diseño.

El aporte metodológico es la medición de cuánto desempeño procede de la propia inspección. En
estos datos, el reemplazo del contador y los cambios de código y de tarifa se concentran al
final del historial de los clientes con fraude, y un modelo que los ve alcanza cifras cercanas
a las de la literatura. Al censurar el historial en el primer evento de ese tipo, la ventaja
de la segunda mitad del historial sobre la primera cae de 0,213 a 0,032 de PR-AUC. La cifra
que queda, un PR-AUC de 0,191 y un ROC-AUC de 0,762 en el conjunto de prueba, es la que cabe
esperar al puntuar clientes aún no inspeccionados, y es la que debería guiar una decisión de
uso.

Los tres objetivos específicos se cumplieron. La caracterización de los datos identificó la
longitud del historial como factor de confusión, con una tasa de fraude que se multiplica por
15,2 entre el decil de clientes con menos facturas y el de más. La comparación de modelos, con
pliegues agrupados por cliente, variables relativas ajustadas dentro de cada pliegue y un
conjunto de prueba evaluado una sola vez, seleccionó LightGBM con hiperparámetros por defecto,
indistinguible de XGBoost y mejor que la regresión logística. La evaluación práctica mostró una
precisión del 43,5\,\% en la cabeza de la lista, probabilidades calibradas y una señal que se
mantiene en los clientes sin eventos de equipo o de contrato.

Para trasladar el método a las distribuidoras del Caribe colombiano, la condición más
importante es de registro: conservar la fecha y el resultado de cada
inspección, también de las que no encuentran fraude. Con esa información se puede separar el
historial anterior a la detección del posterior, validar con un corte temporal y fijar el
tamaño de la lista a partir de los costos reales de cada visita.
